\documentclass[../main.tex]{subfiles}
\begin{document}
%==========================================TIÊU ĐỀ TẬP========================================
\begin{table}[h]
  \begin{adjustwidth}{-1cm}{}
    \begin{tabular}{>{\centering\arraybackslash}p{6cm}|>{\raggedright\arraybackslash}p{12.5cm}}
      \multirow{5}{6cm}
      % Cột bên trái: ÉPISODE và số tập
      {\\[-40pt]
        \flaregothic\fontsize{20pt}{20pt}\selectfont \centering
        \color{eptype}ÉPISODE\color{black}\\
        \burbank\fontsize{72pt}{72pt}\selectfont
        \color{epnum}20\color{black}\\[6pt]
        \cafeteria\fontsize{20pt}{20pt}\selectfont\color{epnum}(2-3)\color{black}
        \\[18pt]
        \flaregothic\fontsize{18pt}{18pt}\selectfont
        \color{epattr}ÉPISODE PERDU
        \color{black}
      }
       &
      % Dòng thứ nhất: tên tập tiếng Nhật
      {\fotskip\fontsize{18pt}{18pt}\selectfont \ruby{新}{しん}\ruby{衣}{い}\ruby{装}{しょう}お\ruby{披}{ひ}\ruby{露}{ろ}\ruby{目}{め}！？
      } \\[6pt]
      % Dòng thứ hai: tên tập tiếng Anh
       & {\cafeteria\fontsize{18pt}{18pt}\selectfont The Twin-tail Heist}                                               \\[4pt]
      % Dòng thứ ba: tên tập tiếng Việt
       & {\vnmsans\fontsize{18pt}{18pt}\selectfont Kẻ cắp đôi bím tóc}                                                  \\[8pt]
      % Dòng thứ tư: Nhân vật xuất hiện
       & {\caviar{\fontsize{14pt}{14pt}\selectfont Nhân vật: \emp{kuon1} \emp{ryuuhou1} ({\abv Truyện phụ}) \emp{miko} \emp{aki} \emp{mio}}} \\[8pt]
      % Dòng cuối cùng: Thời gian diễn ra của tập (thường sẽ là ngày phát hành)
       & {\continuum\fontsize{16pt}{16pt}\selectfont Ngày phát hành: 22/09/2019}                                        \\[18pt]
    \end{tabular}
  \end{adjustwidth}
\end{table}
%=========================================TRUYỆN CHÍNH========================================
\color{black}

\txtbx[2]{{\color{vol} \uline{Tóm tắt:}} Đôi bím tóc của Aki bất ngờ ``vượt ngục'' và bay lượn tự do giữa bầu trời, để lại chủ nhân trong sự hoảng loạn. Khi nỗ lực quay trở về, chúng đã gây ra một vụ tai nạn dở khóc dở cười cho Aki và Mio, đồng thời trở thành món đồ chơi mới cho buổi livestream đầy chiêu trò của Miko.}

\pov{Hikawa Kuon}

Một buổi sáng nọ, tôi cùng nàng Sói đen Mio đang ngồi đối diện nhau trong phòng hội nghị để bàn bạc về ý tưởng cho buổi công chiếu sắp tới. Trên bàn là những chồng giấy tờ lộn xộn, laptop đang mở với tiếng gõ phím đều đều của tôi. Không khí khá yên tĩnh, chỉ có tiếng quạt máy chạy rì rì và thỉnh thoảng là vài câu trao đổi ngắn giữa hai chúng tôi.

Đột nhiên, cửa phòng bật mở cái rầm khiến cả hai giật mình quay lại. Aki lao vào với đôi mắt xanh lấp lánh ánh hoảng loạn. Giọng em ấy gần như nghẹn lại: ``\textbf{Chị\dots\ chị làm mất đôi bím tóc rồi!}''

\img{2-4a}

Lời nói thất thanh của em ấy khiến chúng tôi ngơ ngác. Quả nhiên, cô nàng Dị giới chỉ còn lại phần tóc trên đầu, còn hai bím tóc vốn lơ lửng đằng sau đã biến mất không dấu vết. 

Mio nhướng mày, đôi tai sói khẽ động đậy vì ngạc nhiên: ``Làm thế nào mà chị lại đánh mất được?'' 

Tôi thì nheo mắt đầy nghi hoặc, cố giữ cho em ấy bình tĩnh lại: ``Anh tưởng chúng nó có thể tự bay được cơ mà?''

Aki hít một hơi thật sâu, giọng của em ấy bất chợt trở nên ủy mị lạ thường khiến chúng tôi bắt đầu đổ mồ hôi hột: ``Một cơn gió thổi bay chúng mất rồi\dots''

Cả tôi và Mio đều thở dài trước cái lý do lãng xẹt đó. 

Sói đen nhìn Aki với ánh mắt cạn lời: ``\dots Có thế thôi ạ?'' 

Tôi thì khoanh tay, châm chọc thêm một câu: ``Nghe nó lãng xẹt quá nhỉ\dots\ Chắc cơn gió đó phải mạnh như siêu bão luôn thì mới thổi bay được cả tóc của em.''

Bất chợt, nàng Sói đen dựng đôi tai của mình lên như nghe thấy thứ gì đó lạ thường. Em ấy quay phắt ra phía cửa sổ, vẻ mặt đầy căng thẳng: ``Tiếng gì vậy?'' 

Đúng lúc đó, tôi cũng nghe thấy tiếng ù ù từ đằng xa, mỗi lúc một lớn dần, nghe y hệt tiếng động cơ của máy bay phản lực đang lao tới.

\img{2-4b}

Từ phía chân trời, hai bím tóc vàng óng quen thuộc đang bay về phía văn phòng với tốc độ kinh hồn như hai quả tên lửa định vị. Aki reo lên đầy vui mừng, đôi mắt sáng rực: ``Chúng nó đang quay lại với chị này!''

Cả tôi và Mio trố mắt nhìn cảnh tượng đó, thực sự không tin nổi vào mắt mình. 

Mio lẩm bẩm trong sự kinh hãi: ``Công nghệ từ dị giới\dots\ đúng là đáng sợ thật.'' 

Tôi chỉ biết gật đầu tán thành: ``Công nhận, ngoài sức tưởng tượng của anh luôn.''

Tuy nhiên, ngay khi hai bím tóc định lao vào phòng qua cửa sổ, chúng đập thẳng vào lớp kính một tiếng *BỊCH* rõ to rồi rơi tuột xuống men cửa như hai sợi bún thiu.

\img{2-4c}

Aki hét lên thảm thiết: ``\textbf{Bím tóc của chị!}'' 

Cùng lúc đó, điện thoại của Mio rung lên báo tin nhắn mới với tiếng chuông ``\textbf{(L)INE!}'' quen thuộc. Chúng tôi kiểm tra tin nhắn, hóa ra là từ A-chan.

\begin{figure}[H]
  \centering
  \begin{subfigure}[b]{0.45\textwidth}
    \centering
    \includegraphics[width=8cm]{../../../../Image/Book2/2-4d1.png}
    \caption{Aqua làm vỡ (từ {\flaregothic ÉPISODE 0}, quyển số Không).}
  \end{subfigure}
  \quad
  \begin{subfigure}[b]{0.45\textwidth}
    \centering
    \includegraphics[width=8cm]{../../../../Image/Book2/2-4d2.png}
    \caption{Hai con cá mập xoay làm vỡ ({\flaregothic ÉPISODE 14}, quyển số Một).}
  \end{subfigure}
  \quad
  \begin{subfigure}[b]{0.45\textwidth}
    \centering
    \includegraphics[width=8cm]{../../../../Image/Book2/2-4d3.png}
    \caption{Ayame làm vỡ (từ {\flaregothic ÉPISODE 19}, quyển số Hai).}
  \end{subfigure}
\end{figure}

Chị ấy nhắn với chúng tôi bằng giọng điệu cực kỳ tự tin: ``Vì cửa sổ văn phòng bị đập vỡ mãi nên chị đã cho lắp kính chống đạn rồi. Yên tâm mà làm việc nhé!''

Nghe vậy, Aki-chan đứng sững người, lấy hai tay che miệng, mắt rơm rớm nước vì tuyệt vọng: ``K-Không thể nào\dots!''

Tôi nhếch môi cười đắc ý, thong thả khoanh tay: ``Chính anh là người đề xuất ý tưởng này với A-san đấy. Cái gì cần bảo vệ thì phải bảo vệ, công ty không thể cứ suốt ngày đền tiền vặt cho mấy cái vụ phá hoại linh tinh này nữa.''

Aki và Mio nhìn tôi với ánh mắt vừa thán phục vừa có chút kinh hãi: ``Đ-Đúng vậy ha\dots''

*BỤP!* *BỤP!*

Bất chợt, tiếng bập bùng lại vang lên từ bên ngoài. Aki quay ra, thấy hai bím tóc của mình đang dính chặt vào kính, cố gắng đập cửa để được vào trong. Em ấy cuồng nhiệt cổ vũ: ``Chúng chưa bỏ cuộc kìa! Cố gắng lên các em! Cố lên nào!''

\gif{2-4e}{1}{120}

Thấy hai lọn tóc cứ đập kính liên hồi, tôi thở dài ngao ngán, bước tới gần cửa sổ rồi tặc lưỡi buông một câu như thể là lẽ hiển nhiên: ``Mở cái cửa sổ ra là xong mà, có gì đâu phải xoắn.'' 

Mio nghe thế liền gật đầu, nhanh tay mở toang cửa sổ ra.

Aki hốt hoảng hét lên cảnh báo nhưng không kịp: ``Khoan đã, nếu em làm như thế\dots''

*VÚT!* *BỐP!*

``\textbf{HỰ!}''

\img{2-4f}

Một tiếng ``uỵch'' khô khốc vang lên, kéo theo tiếng rên đau đớn khi hai bím tóc lao vào với vận tốc bàn thờ, đâm sầm vào Aki và Mio khiến cả hai ngã lăn quay ra sàn, nằm bất tỉnh nhân sự. Tôi đứng ngay sát đó nhưng may mắn né kịp, chỉ biết đổ mồ hôi hột nhìn hiện trường: ``T-Thật là vãi chưởng\dots\ Ăn trọn đôi luôn ạ!''

Tôi chưa kịp hoàn hồn thì cửa phòng lại bật mở một lần nữa. Miko xuất hiện, đúng lúc chứng kiến cảnh tượng dở khóc dở cười: hai bím tóc đang run bần bật như đang cầu cứu trên người hai nạn nhân đang nằm bất động.

\img{2-4g}

Miko bước vào phòng với vẻ mặt nghiêm nghị như một thanh tra cảnh sát, em ấy cầm điện thoại lên và thản nhiên gọi: ``A-lô, đây có phải là cảnh sát không ạ? Tôi muốn báo cáo một vụ tai nạn nghiêm trọng. À, có Kuon-senpai ở đây làm chứng rồi nhe.''

Tôi chỉ biết lắc đầu trước cái tính hay ``xì-căng-đan hóa'' mọi việc của nàng Vu nữ: ``\textit{Tóc mà cũng làm hung thủ được à? Đúng là chỉ có ở văn phòng này mới thế được\dots}''

Hai bím tóc vàng đột nhiên khẽ di chuyển, rồi phát ra những tiếng chí chóe khẩn thiết như đang muốn thanh minh: 
``Khoan đã!'' bím tóc bên trái kêu lên.
``Chúng tôi vô tội mừ!'' bím tóc bên phải hùa theo.

Tôi trố mắt nhìn, thực sự không ngờ tóc của Aki còn biết nói chuyện. Tôi điềm tĩnh đáp lại chúng: ``Ờm\dots\ Không đâu. Về cơ bản, hai người đã phạm tội vô ý gây thương tích rồi đấy. Cứ chuẩn bị tinh thần đi.''

\img{2-4h}

Miko liếc sang tôi, giọng đầy tính đồng minh: ``Kuon-senpai cũng đã chứng kiến hết mọi thứ rồi nhe?'' 

Tôi chỉ khẽ ``ừ'' một tiếng cho xong chuyện, chẳng muốn dây dưa thêm vào đống rắc rối này.

Thấy có sự ủng hộ của tôi, nàng Vu nữ nở một nụ cười nham hiểm, nhìn thẳng vào hai bím tóc đang run rẩy: ``Vậy các người biết phải làm gì rồi đúng không nhe? Đừng để tôi phải ra tay, hậu quả khó lường lắm đó.''

Hai bím tóc sợ đến mức co rúm lại, không dám cãi thêm một câu nào. Bản thân tôi cũng chẳng biết Miko định làm gì, thôi thì cứ tặc lưỡi cho qua vậy.

\ct

``Ờm\dots\ Miko-chan, em định làm gì với cặp tóc đó vậy?'' tôi hỏi khi thấy em ấy xách hai bím tóc trên tay, khuôn mặt hớn hở thấy rõ.

Miko chỉ cười một cách đầy bí hiểm, vẫy tay chào tôi: ``Rồi anh sẽ thấy! Chào nhe!'' Nói đoạn, em ấy chạy biến khỏi phòng nhanh như một cơn lốc, để tôi đứng đó ngơ ngác giữa căn phòng hỗn loạn.

Nhìn lại Aki và Mio vẫn đang nằm vật vã trên sàn, tôi đành thở dài, nhọc nhằn dìu từng người lên ghế sofa nằm nghỉ. Trước khi rời đi, tôi viết vội một tờ giấy nhắn để lại:

\begin{tcolorbox}[
    colframe=vol,
    colback=white,
    boxrule=0.5pt,
    arc=0mm,
    left=5pt,
    right=5pt,
    top=5pt,
    bottom=5pt
  ]
  Aki-chan, Mio-chan,

  Anh phải về trường để giải quyết việc gấp ở CLB. Hai người cứ nghỉ ngơi cho khỏe nhé. Nhớ uống nhiều nước đấy.

  \hfill\textbf{-Hikawa Kuon.}
\end{tcolorbox}

Tôi dán tờ giấy lên bàn, thu dọn đồ đạc rồi lặng lẽ rời khỏi phòng.

\ct

Đi bộ dọc hành lang, tôi vẫn không thôi thắc mắc về ý đồ của Miko. Nhớ lại ánh mắt đầy quyết tâm của nàng Vu nữ lúc nãy, tôi vừa thấy tò mò, vừa có chút e ngại cho số phận của đôi bím tóc tội nghiệp kia. Nhưng thôi, đó là chuyện của tương lai, còn bây giờ tôi phải lo việc của mình đã.

Tôi bước ra khỏi tòa nhà, để lại phía sau một buổi sáng náo nhiệt theo đúng phong cách của hololive.

%==========================================TRUYỆN PHỤ=========================================
\omk

Sáng hôm sau tại nhà Hikawa, tôi đang ngồi nhâm nhi cốc cà phê, tận hưởng buổi sáng yên bình thì điện thoại rung liên hồi. Tin nhắn từ Miko-chi hiện ra: ``Kuon-senpai, vào stream của em ngay nhe! Có bất ngờ cực lớn cho anh đấy nhe\~{}''

Tò mò, tôi mở livestream của nàng Vu nữ lên. Ngay khi hình ảnh hiện ra, tôi suýt chút nữa đã phun hết ngụm cà phê vào màn hình. Cô nàng đang rạng rỡ xuất hiện, nhưng trên đầu em ấy lại là cặp bím tóc vàng của Aki, được gắn một cách hoàn hảo như thể sinh ra là để dành cho em ấy!

Ryuuhou đang ngồi bên cạnh tôi, tay cầm túi snack nhai nhóp nhép, cũng suýt nữa sặc khi nhìn thấy màn hình. Con bé cười nắc nẻ, chỉ tay vào tivi: ``\vie{Onii-san nhìn kìa! Chị Miko gắn cái gì lên đầu vậy? Nhìn như hai cái đuôi tôm chiên khổng lồ ấy! Ha ha ha!}''

\img{2-4i}

Tôi cũng không nhịn được mà phá lên cười: ``\vie{Miko-chan đúng là\dots\ lại có chiêu trò mới rồi. Em ấy bá đạo thật sự.}''

Dõi mắt xuống phần trò chuyện, tôi thấy cộng đồng mạng cũng đang bùng nổ. Những dòng bình luận chạy nhanh như điện:

\begin{itemize}
  \item Kawaii quá trời luôn! Miko-chi lên đời rồi!
  \item Bộ dạng này hợp đến bất ngờ đấy chứ, nhìn cưng xỉu!
  \item Chịu bà rồi đấy Miko-chi ạ. Nhưng mà dễ thương thật sự.
\end{itemize}

Tôi bật cười khi nhìn thấy phản ứng tích cực tới bất ngờ từ mọi người, nhưng chỉ vài giây sau, bình luận của Aki xuất hiện, hòa chung với những lời khen của những người hâm mộ:

\begin{itemize}
  \item Miko-chan\dots
  \item Trả nó lại cho em\dots
\end{itemize}

Lập tức, hàng loạt bình luận ùa vào ``chọc quê'' nàng tóc vàng đen đủi ấy:

\begin{itemize}
  \item Tội nghiệp Aki-chan quá!
  \item Đi đâu mà rơi tóc thế này! www
  \item Chắc Aki-san phải nhớ lắm đây www
  \item Miko-chan bá đạo thật.
  \item wwwwwwwwwwwwwwwwww
\end{itemize}

Ryuuhou hóng hớt đọc bình luận rồi lại quay sang trêu tôi: ``\vie{Onii-san cũng vào bình luận cái gì đi chứ, đừng có ngồi đó mà cười một mình!}''

Chiều lòng em gái, tôi nhanh tay gõ một dòng:
\begin{itemize}
  \item Miko-chan, anh nghĩ là em nên nhuộm chúng sang màu đỏ cho trùng với màu tóc của mình thì sẽ hợp hơn đấy. Chứ để màu vàng thế này trông cứ như đang đeo đồ đi mượn vậy!
\end{itemize}

Ngay lập tức, dòng bình luận của tôi bị fan hâm mộ ``tấn công'' bằng hàng loạt sự đồng tình:
\begin{itemize}
  \item Kuon-san nói chuẩn quá! www
  \item Tóc đỏ đi chị ơi, tạo thương hiệu mới luôn!
\end{itemize}

Trên màn hình, Miko bỗng chốc xanh mặt lại, chắc là em ấy bắt đầu cảm thấy có gì đó ``sai sai'' khi bị tôi bóc phốt ngay trên sóng. Hai anh em tôi ngồi trước laptop cười đến chảy cả nước mắt. 

Nhưng đúng như người ta nói, nụ cười của người này lại là nỗi đau của người khác.

\ct

\img{2-4j}

Tại văn phòng \hll, Aki đã sớm biết chuyện nhờ những đoạn cắt được lan truyền với tốc độ chóng mặt. Nàng tóc vàng phồng má, mặt đỏ bừng vì vừa giận vừa xấu hổ khi thấy tóc mình bị đem ra làm trò đùa. 

Giọng em ấy đầy vẻ ấm ức kêu gào: ``Miko-chan dám `mượn' tóc của mình mà không thèm hỏi một tiếng\dots\ Còn anh nữa, Kuon-senpai! Anh còn cổ xúy cho em ấy nhuộm đỏ nữa là sao!''

Đứng phía sau, Mio vừa ôm bụng cười đến chảy nước mắt vừa cố gắng vỗ vai trấn an bạn: ``Thôi nào Aki-chan, trông Miko-chan cũng dễ thương mà\~{} Coi như là đổi mới phong cách đi.''

Lúc tôi vừa bước vào văn phòng, nghe thấy tiếng than vãn của Aki chỉ biết xoa đầu cười gượng gạo: ``Anh chỉ đóng góp ý kiến mang tính xây dựng thôi mà\dots\ Nhưng thôi được rồi, để anh tìm cách xử lý vụ này cho.''

\ct

Buổi chiều hôm đó, tôi hẹn cả ba người gặp nhau ở quán cà phê quen thuộc. Sau một hồi thuyết phục, Miko cũng chịu trả lại tóc cho Aki, nhưng vẫn không quên để lại một câu đầy tiếc nuối: ``Tiếc quá nhe, nếu nhuộm đỏ chắc chắn sẽ thành trend mới cho xem.''

Aki nhận lại cặp bím tóc, nửa mừng nửa lo cảnh báo: ``Lần tới em sẽ không để Miko-chan lấy đi dễ dàng thế đâu nhe!''

Kết thúc một ngày hỗn loạn bằng một bữa tối vui vẻ do tôi móc hầu bao chiêu đãi. Nhìn họ cười đùa bên nhau, tôi khẽ mỉm cười tự nhủ: ``\textit{Dù có rắc rối thế nào, ở bên họ thực sự không bao giờ thấy nhàm chán.}''

``\textit{Có lẽ mình bắt đầu thực sự thích cái văn phòng kỳ lạ này rồi đấy.}''

\end{document}
